\section*{Capítulo \ref{ch:b3:inorganic-materials} --- Materiales inorgánicos y nanomateriales}

\begin{solution}{exo:b3:inorganic-materials:1}
$x^2 \propto t$: un espesor tres veces mayor tarda nueve veces más, \qty{90}{h}.
\end{solution}

\begin{solution}{exo:b3:inorganic-materials:2}
Hidrólisis: $\ce{Si(OC2H5)4 + 4H2O -> Si(OH)4 + 4C2H5OH}$.\\
Condensación: $\ce{2Si(OH)4 -> Si2O(OH)6 + H2O}$.\\
Global: $\ce{Si(OC2H5)4 + 2H2O -> SiO2 + 4C2H5OH}$.
\end{solution}

\begin{solution}{exo:b3:inorganic-materials:3}
Formadores: $\ce{SiO2}$, $\ce{B2O3}$. Modificadores: $\ce{Na2O}$, $\ce{CaO}$. Cada ion óxido del $\ce{Na2O}$ rompe un puente Si--O--Si
en dos extremos Si--O$^-$, cada uno emparejado con un $\ce{Na+}$: la red se corta, el fundido
fluye con más facilidad y funde a menor temperatura.
\end{solution}

\begin{solution}{exo:b3:inorganic-materials:4}
Doce pentágonos (como todo \omterm{def:b3:inorganic-materials:carbon}{fullereno}). $V = 70$, $E = 3V/2 = 105$ enlaces, $F = 2 - V + E = 37$
caras, luego $37 - 12 = 25$ hexágonos.
\end{solution}

\begin{solution}{exo:b3:inorganic-materials:5}
$\sqrt2\,(r_B + r_O) = 1.4142 \times \qty{200.5}{pm} = \qty{283.5}{pm}$. $\ce{SrTiO3}$: $284/283.5 = 1.00$, la perovskita cúbica ideal.
$\ce{BaTiO3}$: $301/283.5 = 1.06$, titanio demasiado pequeño para su octaedro, descentrado: \omterm{def:b3:inorganic-materials:ferroelectric}{ferroeléctrico}.
$\ce{CaTiO3}$: $274/283.5 = 0.97$, calcio demasiado pequeño para su cavidad: octaedros inclinados, simetría
más baja.
\end{solution}

\begin{solution}{exo:b3:inorganic-materials:6}
Si/Al $= 106/86 = 1.23$. $M = 86 \times 23.0 + 86 \times 27.0 + 106 \times 28.1 + 384 \times 16.0 = \qty{13422.6}{g/mol}$; 43 $\ce{Ca^2+}$ por
fórmula: $43/\qty{13422.6}{g/mol} = \qty{3.20}{mmol/g}$.
\end{solution}

\begin{solution}{exo:b3:inorganic-materials:7}
Distancia entre vecinos más próximos: $d = a/\sqrt2 = \qty{288.4}{pm}$. Un agregado de $n$ capas mide $(2n + 1)d$ de diámetro:
$2n + 1 \approx 5/0.2884 = 17.3$, luego $n = 8$ (\qty{4.9}{nm}). $N(8) = 2057$ átomos, $10 \times 64 + 2 = 642$ en la superficie:
el 31\,\%.
\end{solution}

\begin{solution}{exo:b3:inorganic-materials:8}
$h^2/(8\mu R^2)$ con $\mu = \qty{9.109e-32}{kg}$: \qty{0.94}{eV} a \qty{2.0}{nm} y \qty{0.42}{eV} a \qty{3.0}{nm}. Emisión a
$1.74 + 0.94 = \qty{2.68}{eV}$, \qty{463}{nm} (azul), y a $1.74 + 0.42 = \qty{2.16}{eV}$, \qty{574}{nm} (amarillo): el punto más pequeño
emite la luz más azul.
\end{solution}

\begin{solution}{exo:b3:inorganic-materials:9}
$(10, 10)$: sillón, $n - m = 0$, metálico. $(10, 0)$: zigzag, 10 no es múltiplo de 3, \omterm{def:b3:solid-state:classes}{semiconductor}.
$(9, 0)$: zigzag, metálico. $(12, 8)$: ni lo uno ni lo otro (quiral), $n - m = 4$, \omterm{def:b3:solid-state:classes}{semiconductor}.
\end{solution}

\begin{solution}{exo:b3:inorganic-materials:10}
$M = 4 \times 65.4 + 13 \times 16.0 + 24 \times 12.0 + 12 \times 1.0 = \qty{769.6}{g/mol}$; $a^3 = (\qty{25.82e-8}{cm})^3 = \qty{1.721e-20}{cm^3}$;
$\rho = 8 \times 769.6/(\num{6.022e23} \times \num{1.721e-20}) = \qty{0.594}{g.cm^{-3}}$. Un campo mide \qty{7140}{m^2}: un gramo equivale a
$3000/7140 = 0.42$ campos, y una cucharadita de unos pocos gramos, a más de un campo entero.
\end{solution}

\begin{solution}{exo:b3:inorganic-materials:11}
$\dfrac{10n^2 + 2}{(10n^3 + 15n^2 + 11n + 3)/3} = \dfrac{30n^2 + 6}{10n^3 + 15n^2 + 11n + 3} \to \dfrac{30n^2}{10n^3} = \dfrac3n$. Para $n = 8$: valor exacto
$642/2057 = 0.31$, frente a $3/8 = 0.375$; la aproximación sobrestima en los agregados pequeños,
cuyas capas interiores no son despreciables.
\end{solution}

\begin{solution}{exo:b3:inorganic-materials:12}
$2E = 3V$, $2E = 5p + 6h + 7s$, $F = p + h + s$. Euler: $V - E + F = 2$, con $V = 2E/3$: $F - E/3 = 2$, luego
$p + h + s - (5p + 6h + 7s)/6 = 2$, es decir, $(p - s)/6 = 2$: $p - s = 12$. Cada heptágono exige un
pentágono adicional.
\end{solution}

\begin{solution}{pb:b3:inorganic-materials:1}
\textbf{1.} $\mathrm{Na_{12}Al_{12}Si_{12}O_{48}}$: $12 \times 23.0 + 12 \times 27.0 + 12 \times 28.1 + 48 \times 16.0 = \qty{1705.2}{g/mol}$.

\textbf{2.} $1705.2 + 27 \times 18.0 = \qty{2191.2}{g/mol}$.

\textbf{3.} Si/Al $= 12/12 = 1$.

\textbf{4.} No hay enlaces Al--O--Al; con Si/Al $= 1$, cada Si tiene cuatro vecinos Al y cada
Al cuatro Si: alternancia estricta.

\textbf{5.} $-12$, uno por aluminio.

\textbf{6.} $486.0/2191.2 = 22.2$\,\%.

\textbf{7.} $\ce{2Na+}(\text{zeolita}) + \ce{Ca^2+}(\text{aq}) \rightleftharpoons \ce{Ca^2+}(\text{zeolita}) + \ce{2Na+}(\text{aq})$.

\textbf{8.} Seis (doce cargas).

\textbf{9.} $6/\qty{1705.2}{g/mol} = \qty{3.52}{mmol/g}$.

\textbf{10.} $6/\qty{2191.2}{g/mol} = \qty{2.74}{mmol/g}$.

\textbf{11.} $\qty{3.52}{mmol/g} \times \qty{100.1}{g/mol} = \qty{352}{mg}$ de $\ce{CaCO3}$ por gramo.

\textbf{12.} $\qty{2.0}{mmol/L} \times \qty{15}{L} = \qty{30}{mmol}$.

\textbf{13.} $\qty{30}{mmol}/\qty{2.74}{mmol/g} = \qty{11}{g}$.

\textbf{14.} $\qty{10.9}{g}/0.60 = \qty{18}{g}$.

\textbf{15.} $\qty{60}{mmol}$ de $\ce{Na+}$, $\qty{0.060}{mol} \times \qty{23.0}{g/mol} = \qty{1.4}{g}$.

\textbf{16.} Los fosfatos de las aguas residuales alimentan a las algas de ríos y lagos; las proliferaciones
de algas y su descomposición consumen el oxígeno disuelto (eutrofización).

\textbf{17.} El $\ce{Mg^2+}$, más pequeño, retiene con más fuerza sus moléculas de agua; debe desprenderse de ellas
para atravesar la ventana y llegar a los sitios, lo que es lento a las temperaturas de lavado.

\textbf{18.} El ion desnudo mide unos \qty{0.20}{nm} de diámetro, la mitad de la ventana.

\textbf{19.} El ion hidratado es más grande que la ventana: debe desprenderse de parte de su
capa de agua para pasar, y los oxígenos del armazón ocupan el lugar del
agua.

\textbf{20.} El armazón tiene carga negativa y solo intercambia cationes; un
anión grande de cadena larga es repelido y no puede atravesar la ventana, de modo que se queda
en el agua, donde realiza su trabajo.

\textbf{21.} Deshidratada, la \omterm{def:b3:inorganic-materials:zeolite}{zeolita} capta fuertemente el agua en sus cavidades; solo
entran las moléculas más pequeñas que la ventana, de modo que separa las moléculas por su tamaño:
un tamiz a escala molecular.

\textbf{22.} Los iones $\ce{K+}$, más grandes, se sitúan en las ventanas y las estrechan: solo se admiten
moléculas más pequeñas, como el agua.

\textbf{23.} El isómero \textit{para}, esbelto y lineal, cabe en los canales y difunde deprisa; los
isómeros \textit{orto} y \textit{meta}, más anchos, difunden lentamente, permanecen y se isomerizan.

\textbf{24.} La \omterm{def:b3:inorganic-materials:zeolite}{zeolita} A anhidra puede captar, en teoría, \qty{3.52}{mmol} de $\ce{Ca^2+}$ por gramo.
\end{solution}
